%% =====================================================================
%%  Capítulo 8 -- Controladores PID: sintonización
%%  Migrado al formato nuevo (noviembre de 2026) desde Handout6b_Ctrl.tex.
%%  Correcciones de fondo (ver MIGRACION.md y la caja "Correcciones" de la
%%  sección 8.8): pendiente máxima de un primer orden; PI de Ziegler-Nichols
%%  de malla cerrada (0.45 Kcr, Pcr/1.2); Tyreus-Luyben no es de "un cuarto";
%%  efecto de lambda en síntesis directa; IMC inestable F(1/tau_u) = 1; fila
%%  FOPDT de la tabla IMC; notación tau_i, tau_d -> T_i, T_d; theta -> L.
%% =====================================================================

\capitulo[cap:ps]{Controladores PID: sintonización}
  {La autoridad y el dinero se rodean de calma, y se puede medir el alcance de la influencia de una persona por la intensidad del silencio que la rodea.}
  {Hernán Díaz}

\begin{objetivos}[partes/08-pid-sintonizacion.html]
  \item Usar la terminología del control de procesos y reconocer los elementos de un lazo industrial.
  \item Clasificar un proceso por su respuesta al escalón e identificar un modelo de primer orden con tiempo muerto, con la ganancia normalizada.
  \item Sintonizar un PID con reglas de malla abierta (Ziegler-Nichols, Cohen-Coon, coeficiente de ajustabilidad) y de malla cerrada (Ziegler-Nichols, Tyreus-Luyben), y juzgar el resultado por desempeño \emph{y} por robustez.
  \item Diseñar un controlador por síntesis directa y por modelo interno (IMC), y obtener de ellos un PID con un solo parámetro de ajuste, $\lambda$.
\end{objetivos}

%% ---------------------------------------------------------------------
\section{El lazo de un proceso industrial}
\label{sec:ps-terminos}
%% ---------------------------------------------------------------------

El capítulo~\ref{cap:pn} explicó qué hace cada acción del PID; falta decidir cuánto vale cada una. A eso se le llama \textbf{sintonizar}. El problema no está resuelto en la práctica: un estudio muy citado de Van Overschee y De Moor reportó que cerca del $80$\,\% de los PID industriales estaban mal sintonizados, que un $30$\,\% operaba en manual y que un $25$\,\% conservaba los parámetros de fábrica~\citep{VaDe00}. Y no es por falta de reglas: el compendio de O'Dwyer reúne más de cuatrocientas~\citep{ODw06}. Este capítulo presenta unas pocas, las que más se usan y las que mejor enseñan, empezando por el vocabulario de los procesos, que es donde vive el PID.

\paragraph{Terminología.} La \textbf{variable del proceso} (PV) es la que se quiere controlar: nivel, presión, flujo, temperatura, pH. La \textbf{referencia} o \emph{set point} (SP) es el valor deseado. La \textbf{variable manipulada} (MV) es la que se puede mover para corregir, y la \textbf{salida del controlador} (CO, la $u_c$ del capítulo anterior) es la orden que llega al elemento final. A veces se mide algo distinto de la PV: en un control de nivel que manipula el caudal de entrada, la PV es el nivel y la MV es el caudal. El \textbf{offset} es una desviación permanente entre la PV y la SP, y una \textbf{perturbación de carga} es un cambio indeseado en alguna entrada que afecta la PV. Cuando una persona hace los ajustes, el control es \textbf{manual}; cuando los hace el controlador, \textbf{automático}.

\begin{figure}[ht]
\centering
\includegraphics[width=0.9\linewidth]{cap08/lazo_terminos}
\caption{Lazo de un proceso: controlador, elemento final $\Gv$, proceso $\Gp$, perturbación a través de $G_d$ y medición $G_m$ (sensor y transmisor). El controlador solo ve el producto $\Gv\Gp G_m$.}
\label{fig:ps-lazo}
\end{figure}

\paragraph{Componentes.} La \cref{fig:ps-lazo} muestra los bloques. El \textbf{sensor} o elemento primario es el primero que mide la PV; el \textbf{transductor} convierte una magnitud en otra (una presión en una capacitancia, por ejemplo); un \textbf{convertidor} cambia el tipo de señal (corriente a presión, analógica a digital); el \textbf{transmisor} lleva la medida a una señal estándar. El \textbf{elemento final de control} (válvula, bomba, variador) es el que mueve físicamente la MV, y su \textbf{actuador} lo acciona con energía neumática, hidráulica o eléctrica. Completan el lazo los \textbf{indicadores}, que muestran la medida a un operador, y los \textbf{registradores}, que la guardan.

\paragraph{Señales estándar.} La señal neumática va de $3$ a $15$\,psig; la analógica, de $4$ a $20$\,mA (o de $1$ a $5$\,V), con el cero ``vivo'' en $4$\,mA para distinguir un valor mínimo de un cable roto; y la digital usa protocolos como HART, Fieldbus o Modbus. El extremo inferior del rango se llama LRV y el superior URV.

\begin{ejemplo}[label=ej:ps-4a20]{Un transmisor de 4 a 20 mA}
Un transmisor de temperatura está calibrado entre $95$ y $105\,\grados$C. ¿Qué corriente corresponde a $100\,\grados$C? ¿Y a $98\,\grados$C?
\solucion La relación es lineal: $i=4+16\,\dfrac{T-95}{105-95}$\,mA. Para $100\,\grados$C, $i=12$\,mA, la mitad del rango; para $98\,\grados$C, $i=8.8$\,mA.
\end{ejemplo}

\begin{ideaclave}
\textbf{Lo que se sintoniza es todo el lazo.}\enspace El controlador no distingue la válvula, el proceso y el transmisor: ve un solo sistema, de la CO a la medida. Por eso las pruebas de este capítulo se hacen moviendo la CO y mirando la PV medida, y el ``proceso'' identificado incluye la válvula y el transmisor.
\end{ideaclave}

%% ---------------------------------------------------------------------
\section{Tipos de procesos e identificación}
\label{sec:ps-identificacion}
%% ---------------------------------------------------------------------

La forma más común de conocer la dinámica de un proceso industrial es la \textbf{prueba escalón en lazo abierto} (\cref{fig:ps-prueba}): con el controlador en \emph{manual}, se espera a que la PV esté en estado estacionario, se cambia la CO de golpe y se registra la PV. La respuesta dice dónde estarán las dificultades del control: un tiempo muerto grande, una respuesta inversa, una no linealidad. La misma prueba puede hacerse en lazo cerrado, con el controlador en automático y un escalón de referencia; la usaremos en la \cref{sec:ps-malla-cerrada}.

\begin{figure}[ht]
\centering
\includegraphics[width=0.75\linewidth]{cap08/prueba_abierta}
\caption{Prueba escalón en lazo abierto: el controlador en manual entrega un escalón $\Delta u_c$ y se registra la PV.}
\label{fig:ps-prueba}
\end{figure}

\subsection{Seis tipos de respuesta}

La \cref{fig:ps-tipos} reúne las respuestas típicas.
\begin{description}
  \item[(a) Primer orden.] Sube de inmediato y se acerca exponencialmente al valor final, como un termómetro de vidrio delgado sumergido en agua hirviendo. Se comporta como un filtro pasabajos.
  \item[(b) Orden superior.] La más común: varios almacenamientos en serie (procesos \emph{multicapacidad}) dan una curva en forma de S que arranca con pendiente nula.
  \item[(c) Integrador.] No tiene valor final: la PV crece como una rampa mientras la CO esté desplazada. Aparece donde hay almacenamiento sin autorregulación: nivel, presión en tanques cerrados, concentración en lotes, temperatura en recipientes bien aislados.
  \item[(d) Oscilatorio.] Raro en procesos sin controlar, común en sistemas mecánicos y en lazos ya controlados.
  \item[(e) Tiempo muerto.] Nada pasa durante $L$ segundos: transporte por tuberías o bandas, la ducha del capítulo~\ref{cap:in}, retardos de comunicación. Casi siempre aparece combinado con otro tipo.
  \item[(f) Respuesta inversa.] Arranca en la dirección contraria. Es la marca de un cero en el semiplano derecho (\emph{fase no mínima}). El caso clásico es el nivel de una caldera: más agua fría hace colapsar burbujas de vapor y el nivel primero baja.
\end{description}
Los procesos de primer orden e integradores son los más fáciles de controlar, porque la PV responde de inmediato a la CO; los de tiempo muerto y respuesta inversa, los más difíciles. Los oscilatorios son manejables mientras no se les pida rapidez.

\begin{figure}[ht]
\centering
\includegraphics[width=\linewidth]{cap08/f_tipos}
\caption{Respuesta escalón de los seis tipos de proceso: $1/(4s+1)$, $1/(2s+1)^4$, $1/(s(s+1))$, $1/(s^2+0.4s+1)$, $e^{-6s}/(4s+1)$ y $(1-3s)/((2s+1)(4s+1))$.}
\label{fig:ps-tipos}
\end{figure}

\subsection{Ganancia normalizada}

En los controladores industriales la PV y la CO suelen expresarse en \emph{porcentaje de su rango}, y la ganancia del controlador es adimensional. La ganancia del proceso debe estar en las mismas unidades:
\begin{equation}
  \Kproc=\frac{\Delta y/\text{rango de la PV}}{\Delta u_c/\text{rango de la CO}} .
  \label{eq:ps-ganancia}
\end{equation}

\begin{ejemplo}[label=ej:ps-ganancia]{Ganancia de un proceso de temperatura}
La CO va de $0$ a $100$\,\% y la PV se mide entre $10$ y $70\,\grados$C. La CO pasa de $20$ a $40$\,\% y la PV, de $30$ a $60\,\grados$C. Halle $\Kproc$.
\solucion $\Kproc=\dfrac{30/60}{20/100}=2.5$. En unidades físicas serían $30/20=1.5\,\grados$C/\%; la diferencia es el factor $100/60$ de los rangos.
\end{ejemplo}

\codigorepo{cap08_pid_sintonizacion/C8_Ej1_Identificacion.m}{ganancia}{ejemplo~\ref{ej:ps-ganancia}}

\begin{cuidado}
Si no se dice nada, en este capítulo las señales están normalizadas ($0$--$100$\,\%) y $\Kproc=\Delta y/\Delta u_c$. Mezclar una ganancia física con una regla de sintonización pensada para ganancias normalizadas da un $\Kc$ errado por el cociente de los rangos.
\end{cuidado}

\subsection{Modelo de primer orden con tiempo muerto}

Casi todas las reglas de sintonización resumen el proceso con tres números, los del modelo FOPDT (\emph{first order plus dead time}):
\begin{equation}
  \Gp(s)\approx\frac{\Kproc\,e^{-Ls}}{\tau s+1}.
  \label{eq:ps-fopdt}
\end{equation}
Con la curva de reacción (\cref{fig:ps-identificacion}) se obtienen así:
\begin{receta}[Método de la tangente]
\begin{enumerate}
  \item Ubique el punto de máxima pendiente de la PV y trace la tangente por él.
  \item Donde la tangente corta el valor inicial de la PV, lea $L$ (desde el instante del escalón).
  \item \textbf{Método del 63\,\%:} $\tau$ es el tiempo entre $L$ y el instante en que la PV alcanza el $63.2$\,\% de su cambio total. \textbf{Variante de Ziegler y Nichols:} $\tau$ es el tiempo entre $L$ y el corte de la tangente con el valor final.
  \item $\Kproc=\Delta y/\Delta u_c$ con \eqref{eq:ps-ganancia}.
\end{enumerate}
\end{receta}

\begin{figure}[ht]
\centering
\includegraphics[width=0.85\linewidth]{cap08/f_identificacion}
\caption{Ejemplo~\ref{ej:ps-ident}: tangente de máxima pendiente sobre la respuesta de $e^{-2s}/((8s+1)(2s+1))$. El tiempo muerto aparente ($3.0$\,s) es mayor que el retardo puro ($2$\,s), y $\tau_{ZN}>\tau_{63}$.}
\label{fig:ps-identificacion}
\end{figure}

\begin{ejemplo}[label=ej:ps-ident]{Un segundo orden visto como primer orden}
Identifique un FOPDT para $\Gp(s)=e^{-2s}/((8s+1)(2s+1))$ con los dos métodos.
\solucion La pendiente máxima es $0.079$ por segundo, en $t=5.7$\,s, con la PV en $0.213$; la tangente corta el cero en $L\approx3.0$\,s. El $63.2$\,\% se alcanza en $t=12.3$\,s, así que $\tau_{63}\approx9.3$\,s, mientras que la tangente llega al valor final $\tau_{ZN}\approx12.7$\,s después de $L$. El retardo aparente absorbe el polo rápido ($2$\,s de retardo puro más parte de la constante de $2$\,s), y la variante de Ziegler y Nichols da una constante de tiempo mayor.
\end{ejemplo}

\codigorepo{cap08_pid_sintonizacion/C8_Ej1_Identificacion.m}{tangente}{identificación con la función \cmd{identificarFOPDT}}

Dos observaciones. Primera, en un proceso de orden superior el tiempo muerto identificado suele ser mayor que el retardo puro: incluye el efecto de los polos rápidos. Segunda, la variante de Ziegler y Nichols tiende a sobrestimar $\tau$; las reglas posteriores a la suya usan el método del $63$\,\%.

\begin{cuidado}[La pendiente máxima de un primer orden]
En un proceso de primer orden, $y(t)=\Kproc\bigl(1-e^{-(t-L)/\tau}\bigr)$ tiene pendiente $\frac{\Kproc}{\tau}e^{-(t-L)/\tau}$, que es \emph{máxima al comienzo}, en $t=L$, y no en el $63$\,\% (allí vale solo el $37$\,\% del máximo). Por eso en un primer orden puro la tangente arranca justo en $L$ y los dos métodos dan el mismo $\tau$.
\end{cuidado}

\begin{pausa}
Con un primer orden puro sin retardo, $1/(5s+1)$, ¿qué $L$ y qué $\tau$ dan la tangente de Ziegler y Nichols y el método del 63\,\%?
\respuesta{$L=0$ en los dos. La tangente en $t=0$ tiene pendiente $1/5$ y llega a $1$ en $t=5$: $\tau_{ZN}=5$\,s. El $63.2$\,\% se alcanza en $t=5$\,s: $\tau_{63}=5$\,s.}
\end{pausa}

\subsection{Procesos integradores}

En un integrador lo difícil es empezar: hay que encontrar una CO que deje la PV quieta (el caudal de salida igual al de entrada, por ejemplo). Desde ahí, un escalón $\Delta u_c$ produce, tras un tiempo muerto, una rampa de pendiente $v$ (\cref{fig:ps-integrador}). Como la ganancia estática es infinita, se usa la \textbf{ganancia de velocidad}
\begin{equation}
  K_v=\frac{v}{\Delta u_c},\qquad \Gp(s)\approx\frac{K_v\,e^{-Ls}}{s},
  \label{eq:ps-kv}
\end{equation}
con $L$ medido donde la asíntota de la rampa corta el nivel inicial.

\begin{figure}[ht]
\centering
\includegraphics[width=0.6\linewidth]{cap08/f_integrador}
\caption{Respuesta de un proceso integrador con tiempo muerto: rampa de pendiente $v$ tras un tiempo muerto $L$.}
\label{fig:ps-integrador}
\end{figure}

%% ---------------------------------------------------------------------
\section{Reglas de malla abierta}
\label{sec:ps-malla-abierta}
%% ---------------------------------------------------------------------

Un buen PID debe dar un lazo estable y robusto, rechazar las perturbaciones, seguir la referencia con rapidez y suavidad, y no dejar error en estado estacionario. Las reglas de malla abierta buscan eso a partir del modelo \eqref{eq:ps-fopdt}: se hace la prueba escalón, se identifican $\Kproc$, $L$ y $\tau$, y una tabla da $\Kc$, $\Ti$ y $\Td$ del PID ideal $\Kc\bigl(1+\frac{1}{\Ti s}+\Td s\bigr)$. Todas dependen de la \textbf{razón de controlabilidad} $L/\tau$: cuanto mayor, más difícil el proceso.

\subsection{Ziegler y Nichols}

\begin{historia}[Ziegler y Nichols, 1942]
John G. Ziegler y Nathaniel B. Nichols, ingenieros de la Taylor Instruments, publicaron en 1942 dos métodos para ajustar los controladores que vendía su empresa: uno a partir de la curva de reacción en lazo abierto y otro a partir de la oscilación sostenida en lazo cerrado~\citep{ZiNi42}. Antes hubo reglas sueltas (Callender y colaboradores, 1936), pero fueron las suyas las que se volvieron estándar. Su criterio fue empírico: que la respuesta a una perturbación decayera con una razón de un cuarto entre picos sucesivos.
\end{historia}

La regla supone una curva de reacción en forma de S, es decir, un proceso sin integradores ni polos dominantes complejos; con otra forma de respuesta no debe usarse. Mide $\tau$ con la tangente hasta el $100$\,\% (\cref{sec:ps-identificacion}).

\begin{table}[ht]
\centering
\caption{Ziegler-Nichols de malla abierta para $\Kproc e^{-Ls}/(\tau s+1)$ y para el integrador $K_v e^{-Ls}/s$. La fila serie es la misma sintonía escrita en forma serie (capítulo~\ref{cap:pn}).}
\label{tab:ps-zn-abierta}
\renewcommand{\arraystretch}{1.35}
\begin{tabular}{@{}lcccccc@{}}
\toprule
 & \multicolumn{3}{c}{FOPDT} & \multicolumn{3}{c}{Integrador}\\
\cmidrule(lr){2-4}\cmidrule(l){5-7}
 & $\Kc$ & $\Ti$ & $\Td$ & $\Kc$ & $\Ti$ & $\Td$\\
\midrule
P & $\dfrac{\tau}{\Kproc L}$ & --- & --- & $\dfrac{1}{K_vL}$ & --- & ---\\
PI & $\dfrac{0.9\tau}{\Kproc L}$ & $3.33L$ & --- & $\dfrac{0.9}{K_vL}$ & $3L$ & ---\\
PID ideal & $\dfrac{1.2\tau}{\Kproc L}$ & $2L$ & $0.5L$ & $\dfrac{1.2}{K_vL}$ & $2L$ & $0.5L$\\
PID serie & $\dfrac{0.6\tau}{\Kproc L}$ & $L$ & $L$ & $\dfrac{0.6}{K_vL}$ & $L$ & $L$\\
\bottomrule
\end{tabular}
\end{table}

Note que $\Kc$ crece como $\tau/L$: con un tiempo muerto pequeño frente a la constante de tiempo, la regla pide ganancias enormes y tiempos integrales cortos, y el lazo oscila. Por eso no se recomienda para $L/\tau$ pequeño.

\subsection{Cohen y Coon}

Cohen y Coon partieron del mismo modelo, pero ajustaron las fórmulas para ubicar los polos dominantes del lazo con razón de decaimiento de un cuarto en un rango amplio de $L/\tau$~\citep{CoCo53}. Geraldine Coon verificó el método con simulaciones en un computador analógico. Una ventaja es que incluye un PD. Con $r=L/\tau$:

\begin{table}[ht]
\centering
\caption{Cohen-Coon, con $r=L/\tau$.}
\label{tab:ps-cc}
\renewcommand{\arraystretch}{1.5}
\begin{tabular}{@{}lccc@{}}
\toprule
 & $\Kc$ & $\Ti$ & $\Td$\\
\midrule
P & $\dfrac{\tau}{\Kproc L}\Bigl(1+\dfrac{r}{3}\Bigr)$ & --- & ---\\
PI & $\dfrac{\tau}{\Kproc L}\Bigl(0.9+\dfrac{r}{12}\Bigr)$ & $L\,\dfrac{30+3r}{9+20r}$ & ---\\
PD & $\dfrac{\tau}{\Kproc L}\Bigl(\dfrac54+\dfrac{r}{6}\Bigr)$ & --- & $L\,\dfrac{6-2r}{22+3r}$\\
PID & $\dfrac{\tau}{\Kproc L}\Bigl(\dfrac43+\dfrac{r}{4}\Bigr)$ & $L\,\dfrac{32+6r}{13+8r}$ & $\dfrac{4L}{11+2r}$\\
\bottomrule
\end{tabular}
\end{table}

El precio de ser tan agresivo es la fragilidad: un cambio pequeño en el proceso puede llevar el lazo a la inestabilidad. La \cref{sec:ps-robustez} lo cuantifica.

\subsection{Coeficiente de ajustabilidad}

Como Ziegler-Nichols y Cohen-Coon dan respuestas muy oscilatorias, en la práctica se usan versiones más suaves. Una regla empírica de instrumentistas usa el \textbf{coeficiente de ajustabilidad} $r_c=L/\tau$ (\cref{tab:ps-ajust}); se recomienda para $0.1<r_c<1$, y para $r_c$ grandes un PID solo no basta. La tabla no tiene una derivación publicada que la respalde: tómela como una regla de taller que da resultados conservadores, no como un resultado.

\begin{table}[H]
\centering
\caption{Regla del coeficiente de ajustabilidad, $r_c=L/\tau$.}
\label{tab:ps-ajust}
\renewcommand{\arraystretch}{1.5}
\begin{tabular}{@{}lccc@{}}
\toprule
$r_c$ & $\Kc$ & $\Ti$ & $\Td$\\
\midrule
$0$ a $0.1$ & $5/\Kproc$ & $\tau$ & $0$\\
$0.1$ a $0.2$ & $\dfrac{0.5}{\Kproc r_c}$ & $\tau$ & $0$\\
$0.2$ a $0.5$ & $\dfrac{0.5(1+0.5r_c)}{\Kproc r_c}$ & $\tau(1+0.5r_c)$ & $\tau\,\dfrac{0.5r_c}{0.5r_c+1}$\\
\bottomrule
\end{tabular}
\end{table}

\begin{ejemplo}[label=ej:ps-abierta]{Tres reglas sobre el mismo proceso}
Un proceso se identificó como $2e^{-2s}/(10s+1)$. Sintonice un PID ideal con las tres reglas y compare la respuesta a un escalón de referencia y a un escalón de perturbación a la entrada del proceso. Use la estructura PI-D (derivada solo sobre la medida) y filtro derivativo con $N=10$.
\solucion Con $L/\tau=0.2$:
\begin{center}
\renewcommand{\arraystretch}{1.2}
\begin{tabular}{@{}lccccccc@{}}
\toprule
Regla & $\Kc$ & $\Ti$ (s) & $\Td$ (s) & $\OS$ & $\ts$ (s) & pico pert. & IAE pert.\\
\midrule
Ziegler-Nichols & $3$ & $4$ & $1$ & $64$\,\% & $19.7$ & $0.38$ & $1.6$\\
Cohen-Coon & $3.46$ & $4.55$ & $0.70$ & $84$\,\% & $24.2$ & $0.39$ & $1.6$\\
Ajustabilidad & $1.38$ & $11$ & $0.91$ & $2$\,\% & $20.4$ & $0.48$ & $8.0$\\
\bottomrule
\end{tabular}
\end{center}
Ziegler-Nichols y Cohen-Coon responden rápido pero con sobrepicos enormes; la regla de ajustabilidad casi no sobrepasa, pero deja crecer más la perturbación y tarda mucho más en eliminarla: su IAE ante la perturbación es cinco veces mayor, porque $\Ti$ es casi tres veces más largo (\cref{fig:ps-abierta}). Ante la perturbación, Ziegler-Nichols decae con una razón cercana a $0.19$ entre picos sucesivos: cerca del cuarto que buscaba.
\end{ejemplo}

\begin{figure}[ht]
\centering
\includegraphics[width=\linewidth]{cap08/f_malla_abierta}
\caption{Ejemplo~\ref{ej:ps-abierta}: tres sintonías de malla abierta sobre $2e^{-2s}/(10s+1)$.}
\label{fig:ps-abierta}
\end{figure}

\codigorepo{cap08_pid_sintonizacion/C8_Ej2_MallaAbierta.m}{reglas}{las tres reglas con la función \cmd{reglasPID}}
\codigorepo{cap08_pid_sintonizacion/C8_Ej2_MallaAbierta.m}{simulacion}{referencia y perturbación con \cmd{simularPID}}

\begin{ideaclave}
\textbf{Seguir la referencia o rechazar perturbaciones.}\enspace Ziegler-Nichols y Cohen-Coon se diseñaron para rechazar perturbaciones, y en eso son buenos; ante cambios de referencia sobrepasan demasiado. Una regla conservadora hace lo contrario. No hay una sintonía mejor en todo: hay que decidir para qué se sintoniza.
\end{ideaclave}

\begin{pausa}
En el ejemplo~\ref{ej:ps-abierta}, ¿qué pasa con el $\Kc$ de Ziegler-Nichols si el tiempo muerto fuera de $0.5$\,s en lugar de $2$\,s?
\respuesta{Se multiplica por cuatro: $\Kc=1.2\cdot10/(2\cdot0.5)=12$, y $\Ti=1$\,s. Un $L/\tau$ pequeño lleva a ganancias muy altas.}
\end{pausa}

%% ---------------------------------------------------------------------
\section{Reglas de malla cerrada}
\label{sec:ps-malla-cerrada}
%% ---------------------------------------------------------------------

Las reglas de malla abierta exigen sacar el lazo a manual, y eso no siempre es posible: un proceso inestable en lazo abierto (un reactor exotérmico, ciertos niveles) se escapa durante la prueba. Las de malla cerrada dejan el controlador en \emph{automático} y caracterizan el proceso con dos números medidos en el propio lazo.

\subsection{Ajuste manual}

Antes de las reglas, el método que se usa en campo: ensayo y error con una secuencia ordenada. Es lento y la acción derivativa es difícil de ajustar así, pero no necesita modelo y se inspira en la misma idea que la regla de Ziegler y Nichols.
\begin{receta}[Ajuste manual en lazo cerrado]
\begin{enumerate}
  \item Arranque con $\Kc$ y $\Td$ pequeños y $\Ti$ grande.
  \item Duplique $\Kc$ y observe la respuesta, hasta que aparezcan oscilaciones sostenidas con $\Kc=K_{cu}$. Deje $\Kc=K_{cu}/2$.
  \item Reduzca $\Ti$ a la mitad sucesivamente hasta que vuelva a oscilar, con $\Ti=\Ti^*$. Deje $\Ti=2\Ti^*$.
  \item Aumente $\Td$ hasta que la respuesta oscile, con $\Td=\Td^*$. Deje $\Td=\Td^*/3$.
\end{enumerate}
\end{receta}

\subsection{Ziegler y Nichols de malla cerrada}

Con el controlador en solo proporcional ($\Ti\to\infty$, $\Td=0$), se sube $\Kc$ desde cero hasta que la salida oscila sin crecer ni decaer. Esa es la \textbf{ganancia crítica} $\Kcr$, y el periodo de la oscilación es el \textbf{periodo crítico} $\Pcr$. En ese punto dos polos del lazo cerrado están sobre el eje imaginario, en $\pm j\wcr$ con $\wcr=2\pi/\Pcr$. Si el lazo nunca oscila al subir la ganancia (un primer o segundo orden sin retardo, por ejemplo), el método no se aplica. Con un modelo, $\Kcr$ y $\wcr$ salen de sustituir $s=\jw$ en la ecuación característica, o de la fila nula de Routh (capítulo~\ref{cap:es}).

\begin{table}[ht]
\centering
\caption{Reglas de malla cerrada: Ziegler-Nichols~\citep{ZiNi42} y Tyreus-Luyben~\citep{TyLu92}.}
\label{tab:ps-cerrada}
\renewcommand{\arraystretch}{1.3}
\begin{tabular}{@{}llccc@{}}
\toprule
Regla & & $\Kc$ & $\Ti$ & $\Td$\\
\midrule
Ziegler-Nichols & P & $0.5\Kcr$ & --- & ---\\
 & PI & $0.45\Kcr$ & $\Pcr/1.2$ & ---\\
 & PID ideal & $0.6\Kcr$ & $0.5\Pcr$ & $0.125\Pcr$\\
 & PID serie & $0.3\Kcr$ & $0.25\Pcr$ & $0.25\Pcr$\\
\midrule
Tyreus-Luyben & PI & $\Kcr/3.2$ & $2.2\Pcr$ & ---\\
 & PID ideal & $\Kcr/2.2$ & $2.2\Pcr$ & $\Pcr/6.3$\\
\bottomrule
\end{tabular}
\end{table}

\begin{ejemplo}[label=ej:ps-critico]{Ganancia y periodo críticos}
Halle $\Kcr$ y $\Pcr$ para $\Gp(s)=1/((s+1)(0.5s+1)(0.2s+1))$ con control proporcional.
\solucion La ecuación característica es $0.1s^3+0.8s^2+1.7s+1+\Kc=0$, y Routh da estabilidad para $-1<\Kc<12.6$. Con $s=\jw$:
\[
  \underbrace{-0.8\omega^2+1+\Kc}_{\text{real}}+j\,\underbrace{\bigl(-0.1\omega^3+1.7\omega\bigr)}_{\text{imaginaria}}=0 .
\]
La parte imaginaria da $\wcr=\sqrt{17}=4.12$\,rad/s, y la real, $\Kcr=0.8\cdot17-1=12.6$. Entonces $\Pcr=2\pi/\sqrt{17}=1.52$\,s.
\end{ejemplo}

La meta de Ziegler y Nichols era la \textbf{razón de decaimiento de un cuarto}: que cada pico de la oscilación fuera la cuarta parte del anterior. Es una respuesta bastante oscilatoria, con sobrepicos de $40$ a $60$\,\% ante la referencia.

\subsection{Tyreus y Luyben}

Tyreus y Luyben propusieron sus reglas para procesos con integrador y tiempo muerto, donde Ziegler-Nichols resulta demasiado agresivo~\citep{TyLu92}. Usan los mismos $\Kcr$ y $\Pcr$, pero con menos ganancia y tiempos integrales mucho más largos: dan respuestas más amortiguadas y menos sensibles a cambios en el proceso. \alerta{No buscan la razón de un cuarto}: son deliberadamente más conservadoras.

\begin{ejemplo}[label=ej:ps-zntl]{Ziegler-Nichols contra Tyreus-Luyben}
Para $\Gp(s)=1/((3s+1)(2s+1)(s+1))$, halle $\Kcr$ y $\Pcr$, sintonice un PID ideal con las dos reglas y compare.
\solucion La ecuación característica es $6s^3+11s^2+6s+1+\Kc=0$. La parte imaginaria, $-6\omega^3+6\omega=0$, da $\wcr=1$\,rad/s y $\Pcr=2\pi=6.28$\,s; la real, $\Kcr=11-1=10$. Las reglas dan
\begin{center}
\renewcommand{\arraystretch}{1.2}
\begin{tabular}{@{}lcccccc@{}}
\toprule
Regla & $\Kc$ & $\Ti$ (s) & $\Td$ (s) & $\OS$ & $\ts$ (s) & decaimiento\\
\midrule
Ziegler-Nichols & $6$ & $3.14$ & $0.785$ & $46.5$\,\% & $16.3$ & $0.17$\\
Tyreus-Luyben & $4.55$ & $13.8$ & $1.00$ & $9.0$\,\% & $30.0$ & sin segundo pico\\
\bottomrule
\end{tabular}
\end{center}
(simulación con filtro derivativo $N=10$). Ziegler-Nichols oscila con una razón de $0.17$, algo mejor que un cuarto; Tyreus-Luyben casi no sobrepasa, a cambio de un establecimiento más lento (\cref{fig:ps-cerrada}). Si la ganancia del proceso sube un $50$\,\% sin retocar el controlador, el sobrepico pasa a $55$\,\% con Ziegler-Nichols y a $22$\,\% con Tyreus-Luyben.
\end{ejemplo}

\begin{figure}[ht]
\centering
\includegraphics[width=\linewidth]{cap08/f_malla_cerrada}
\caption{Ejemplo~\ref{ej:ps-zntl}. Izquierda: lazo proporcional por debajo, en y por encima de $\Kcr$. Derecha: PID ideal sintonizado por las dos reglas.}
\label{fig:ps-cerrada}
\end{figure}

\codigorepo{cap08_pid_sintonizacion/C8_Ej3_MallaCerrada.m}{critico}{$\Kcr$ y $\Pcr$ por sustitución y con \cmd{tablaRouth}}
\codigorepo{cap08_pid_sintonizacion/C8_Ej3_MallaCerrada.m}{zntl}{las dos sintonías y la razón de decaimiento}

\begin{cuidado}[Llevar un lazo real al límite]
La prueba de malla cerrada lleva el proceso al borde de la inestabilidad. En una planta real eso puede ser inaceptable (o peligroso), y el actuador puede saturar antes de que la oscilación sea limpia. Por eso se prefiere la variante con relé de Åström y Hägglund, que produce una oscilación controlada de amplitud pequeña~\citep{AsHa84}.
\end{cuidado}

\begin{pausa}
Para el ejemplo~\ref{ej:ps-zntl}, ¿qué PI daría Ziegler-Nichols de malla cerrada?
\respuesta{$\Kc=0.45\cdot10=4.5$ y $\Ti=2\pi/1.2=5.24$\,s.}
\end{pausa}

%% ---------------------------------------------------------------------
\section{Síntesis directa}
\label{sec:ps-sintesis}
%% ---------------------------------------------------------------------

Las reglas anteriores fijan la estructura (un PID) y ajustan sus tres números. La síntesis directa invierte el razonamiento: se escoge la respuesta de lazo cerrado que se quiere, $T(s)$, y se despeja el controlador que la produce. De $T=\Gc\Gp/(1+\Gc\Gp)$,
\begin{equation}
  \Gc(s)=\frac{1}{\Gp(s)}\,\frac{T(s)}{1-T(s)} ,
  \label{eq:ps-sintesis}
\end{equation}
la \clave{ecuación de síntesis}. El controlador contiene el inverso del modelo: cancela la dinámica del proceso y la reemplaza por la deseada. Nada garantiza que el resultado sea implementable, y lo difícil es escoger $T(s)$. Debe dar error nulo en estado estacionario ($T(0)=1$), ser rápida sin sobrepasar mucho, y ser simple. Las elecciones usuales son
\begin{equation}
  T(s)=\frac{1}{\lambda s+1}
  \qquad\text{o}\qquad
  T(s)=\frac{1}{(\lambda_1 s+1)(\lambda_2 s+1)},
  \label{eq:ps-T-deseada}
\end{equation}
que llevadas a \eqref{eq:ps-sintesis} dan
\[
  \Gc(s)=\frac{1}{\Gp(s)}\,\frac{1}{\lambda s}
  \qquad\text{y}\qquad
  \Gc(s)=\frac{1}{\Gp(s)}\,\frac{1}{s\bigl(\lambda_1\lambda_2 s+\lambda_1+\lambda_2\bigr)} .
\]
En los dos aparece un integrador: por eso el error ante el escalón es nulo.

\begin{ideaclave}
\textbf{$\lambda$ es la constante de tiempo del lazo cerrado.}\enspace Con $T=1/(\lambda s+1)$, el lazo responde como un primer orden de constante $\lambda$: se establece en unos $4\lambda$. Un $\lambda$ \emph{pequeño} da una respuesta \emph{rápida}, a cambio de una ganancia del controlador grande y una acción de control mayor; un $\lambda$ grande da un lazo lento y suave. Es un único parámetro de ajuste con significado físico.
\end{ideaclave}

\paragraph{Proceso estático.} Si $\Gp=\Kproc$, $\Gc=1/(\lambda\Kproc s)$: un controlador integral puro.

\paragraph{Primer orden.} Si $\Gp=\Kproc/(\tau s+1)$,
\[
  \Gc(s)=\frac{\tau s+1}{\Kproc\lambda s}=\frac{\tau}{\Kproc\lambda}\left(1+\frac{1}{\tau s}\right):
  \qquad \Kc=\frac{\tau}{\Kproc\lambda},\quad \Ti=\tau ,
\]
un PI cuyo cero cancela el polo del proceso.

\paragraph{Segundo orden.} Si $\Gp=\Kproc/((\tau_1s+1)(\tau_2s+1))$,
\[
  \Gc(s)=\frac{(\tau_1s+1)(\tau_2s+1)}{\Kproc\lambda s}:
  \qquad \Kc=\frac{\tau_1+\tau_2}{\Kproc\lambda},\quad \Ti=\tau_1+\tau_2,\quad \Td=\frac{\tau_1\tau_2}{\tau_1+\tau_2},
\]
un PID ideal. En la práctica se le agrega el filtro derivativo.

\begin{ejemplo}[label=ej:ps-sintesis]{Un PID por síntesis directa}
Para $\Gp(s)=3/((2s+1)(10s+1))$ y $T(s)=1/(\lambda s+1)$, halle el PID y evalúe $\lambda=1$, $3$ y $5$\,s.
\solucion $\Kc=12/(3\lambda)=4/\lambda$, $\Ti=12$\,s y $\Td=20/12=1.67$\,s. Con filtro $N=10$, la simulación da
\begin{center}
\renewcommand{\arraystretch}{1.2}
\begin{tabular}{@{}cccc@{}}
\toprule
$\lambda$ (s) & $\Kc$ & $\ts$ (s) & $|u_c|_{\max}$\\
\midrule
$1$ & $4$ & $3.9$ & $43.8$\\
$3$ & $1.33$ & $12.1$ & $14.6$\\
$5$ & $0.8$ & $19.8$ & $8.8$\\
\bottomrule
\end{tabular}
\end{center}
El tiempo de establecimiento es casi exactamente $4\lambda$, y la acción de control crece como $1/\lambda$.
\end{ejemplo}

\codigorepo{cap08_pid_sintonizacion/C8_Ej4_SintesisDirecta.m}{pi}{síntesis directa de un PI}
\codigorepo{cap08_pid_sintonizacion/C8_Ej4_SintesisDirecta.m}{pid}{ejemplo~\ref{ej:ps-sintesis}}

\subsection{Procesos de fase no mínima}

Un cero en el semiplano derecho no puede invertirse: se convertiría en un polo inestable del controlador. Lo mismo vale para el tiempo muerto, que es de fase no mínima (su aproximación de Padé de la \cref{sec:es-retardo} tiene un cero en $+2/L$). La salida es dejar ese factor en la respuesta deseada: lo que el proceso no permite cambiar debe aparecer en $T(s)$.

\begin{ejemplo}[label=ej:ps-nomin]{Síntesis directa con un cero inestable}
Para $\Gp(s)=(1-9s)/((15s+1)(3s+1))$, diseñe por síntesis directa.
\solucion Con $T=1/(\lambda s+1)$,
\[
  \Gc(s)=\frac{(15s+1)(3s+1)}{\lambda s(1-9s)},
\]
que tiene un polo en $s=+1/9$: no es implementable. Si se conserva el cero en la respuesta deseada y se sube el orden para mantener el grado relativo, $T(s)=(1-9s)/(\lambda s+1)^2$, entonces
\[
  \Gc(s)=\frac{45s^2+18s+1}{s\bigl(\lambda^2s+2\lambda+9\bigr)}
  =\Kc\left(1+\frac{1}{\Ti s}+\Td s\right)\frac{1}{\tau_f s+1},
\]
un PID ideal en serie con un filtro de primer orden, con
\[
  \Kc=\frac{18}{2\lambda+9},\qquad \Ti=18\text{ s},\qquad \Td=2.5\text{ s},\qquad \tau_f=\frac{\lambda^2}{2\lambda+9}.
\]
Para $\lambda=5$: $\Kc=0.947$ y $\tau_f=1.32$\,s. La respuesta inversa sigue ahí (\cref{fig:ps-nomin}): ningún controlador lineal estable la elimina.
\end{ejemplo}

\begin{figure}[ht]
\centering
\includegraphics[width=\linewidth]{cap08/f_sintesis_nomin}
\caption{Ejemplo~\ref{ej:ps-nomin}: salida y acción de control para tres valores de $\lambda$. Un $\lambda$ menor acelera la respuesta pero profundiza el arranque en sentido contrario y exige más acción de control.}
\label{fig:ps-nomin}
\end{figure}

\codigorepo{cap08_pid_sintonizacion/C8_Ej4_SintesisDirecta.m}{nomin}{ejemplo~\ref{ej:ps-nomin}}

\begin{pausa}
¿Por qué, en el ejemplo~\ref{ej:ps-nomin}, la respuesta deseada tiene denominador de segundo orden y no de primero?
\respuesta{Para conservar el grado relativo. $\Gp$ tiene grado relativo $1$, y un lazo cerrado no puede tener menos. Con $T=(1-9s)/(\lambda s+1)$, de grado relativo $0$, saldría $\Gc=(15s+1)(3s+1)/((\lambda+9)s)$: impropio, un PID ideal sin filtro que pediría a la salida saltar instantáneamente.}
\end{pausa}

%% ---------------------------------------------------------------------
\section{Control por modelo interno}
\label{sec:ps-imc}
%% ---------------------------------------------------------------------

La síntesis directa usa un modelo del proceso para diseñar. El \clave{control por modelo interno} (IMC, \emph{internal model control}) va más lejos: mete el modelo \emph{dentro} del controlador, y con eso maneja de forma explícita los errores de modelado y las perturbaciones~\citep{GaMo82,RMS86}. Para ciertos modelos, el resultado es un PID con un solo parámetro de ajuste.

\subsection{Control basado en modelo, en lazo abierto}

Supongamos un proceso estable $\Gp$ y un controlador $Q(s)$ en lazo abierto: $Y=Q\Gp R$. Con $\Gp=\Kproc/(\tau s+1)$, el controlador estático $Q=1/\Kproc$ lleva la salida al valor de la referencia, pero con la dinámica del proceso. Para acelerar habría que invertir el proceso, $Q=1/\Gp$, y la salida seguiría a la referencia de manera perfecta. Pero
\[
  Q(s)=\frac{\tau s+1}{\Kproc}
  \quad\Longrightarrow\quad
  U(s)=\frac{1}{\Kproc}(\tau s+1)R(s),
  \qquad
  u(t)=\frac{\tau}{\Kproc}\frac{\dif r}{\dif t}+\frac{r(t)}{\Kproc},
\]
que es impropio: ante un escalón de referencia de tamaño $A$ pide un impulso $\frac{A\tau}{\Kproc}\impulso(t)$. La solución es un filtro que haga $Q$ propio,
\begin{equation}
  F(s)=\frac{1}{\lambda s+1},\qquad Q(s)=\frac{F(s)}{\Gp(s)}=\frac{1}{\Kproc}\,\frac{\tau s+1}{\lambda s+1},
  \label{eq:ps-q-filtro}
\end{equation}
un adelanto-atraso. Ahora $Y=F R$: un primer orden de constante $\lambda$, más rápido que el proceso si $\lambda<\tau$. Como en síntesis directa, $\lambda$ es la constante de tiempo deseada.

\subsection{Estabilidad interna}

¿Se puede invertir cualquier proceso? Con $\Gp=(1-9s)/((15s+1)(3s+1))$ y $\lambda=5$, compare
\[
  Q_1(s)=\frac{(15s+1)(3s+1)}{(1-9s)(5s+1)}
  \qquad\text{y}\qquad
  Q_2(s)=\frac{(15s+1)(3s+1)}{(5s+1)^2} .
\]
$Q_1$ invierte todo el proceso y $\Gp Q_1=1/(5s+1)$: una salida impecable. Pero $Q_1$ tiene un polo en $s=+1/9$, y la acción de control crece como $e^{t/9}$: supera $50$ unidades en $t=28$\,s y diverge (\cref{fig:ps-interna}). La salida no lo muestra porque el cero del proceso cancela exactamente ese polo. $Q_2$ deja el cero sin invertir; la salida conserva la respuesta inversa, $Y=(1-9s)/(5s+1)^2\,R$, pero todas las señales están acotadas.

\begin{figure}[ht]
\centering
\includegraphics[width=\linewidth]{cap08/f_imc_interna}
\caption{Estabilidad interna: si $Q$ invierte el cero en $+1/9$, la salida es perfecta mientras la acción de control diverge.}
\label{fig:ps-interna}
\end{figure}

\begin{definicion}[label=def:ps-interna]{Estabilidad interna}
Un sistema de control es \textbf{internamente estable} si toda señal del lazo (salida, acción de control, error, señales internas) permanece acotada ante cualquier entrada acotada aplicada en cualquier punto. Revisar solo $Y/R$ no basta: una cancelación entre un polo inestable y un cero inestable la esconde.
\end{definicion}

Por eso el modelo se factoriza en una parte no invertible y una invertible,
\begin{equation}
  \tilde G_p(s)=\tilde G_p^+(s)\,\tilde G_p^-(s),
  \label{eq:ps-factorizacion}
\end{equation}
donde $\tilde G_p^+$ contiene los ceros en el semiplano derecho y el tiempo muerto, con $\tilde G_p^+(0)=1$. Hay dos formas: la \textbf{simple}, que deja en $\tilde G_p^+$ solo el factor $(1-as)$; y la \textbf{pasa-todo} (\emph{all-pass}), que lo divide por su espejo, $\tilde G_p^+=(1-as)/(1+as)$, de modo que $|\tilde G_p^+(\jw)|=1$ y $\tilde G_p^-$ gana un polo estable. Solo $\tilde G_p^-$ se invierte:
\begin{equation}
  Q(s)=\bigl(\tilde G_p^-(s)\bigr)^{-1}F(s),
  \label{eq:ps-q}
\end{equation}
con el orden del filtro escogido para que $Q$ sea propio.

\subsection{La estructura IMC}

En lazo abierto, cualquier error del modelo pasa directo a la salida. Si el proceso real es $2/(8s+1)$ y se diseña con el modelo $1/(10s+1)$, $Q=(10s+1)/(\lambda s+1)$ y la salida se va a $2$ en lugar de $1$. Y una perturbación no se corrige nunca. La estructura IMC (\cref{fig:ps-imc}) corre el modelo $\tilde G_p$ en paralelo con el proceso y retroalimenta solo la diferencia, $\tilde D=Y-\tilde G_pU$: lo que el modelo no explica, sea perturbación o error de modelado. De la figura,
\begin{equation}
  Y=\frac{\Gp Q}{1+Q(\Gp-\tilde G_p)}\,R+\frac{1-\tilde G_pQ}{1+Q(\Gp-\tilde G_p)}\,D .
  \label{eq:ps-imc-lazo}
\end{equation}
Con un modelo perfecto, $\Gp=\tilde G_p$, el denominador es $1$: $Y=\Gp QR+(1-\Gp Q)D$. El lazo es estable si $Q$ y $\Gp$ lo son, y $Q(0)=1/\tilde G_p(0)$ asegura que la perturbación constante se rechaza por completo. Con modelo imperfecto, la estabilidad depende de $1+Q(\Gp-\tilde G_p)=0$, y $\lambda$ se vuelve una perilla de robustez.

\begin{figure}[ht]
\centering
\includegraphics[width=0.85\linewidth]{cap08/imc}
\caption{Estructura IMC. Solo vuelve a la entrada lo que el modelo no explica.}
\label{fig:ps-imc}
\end{figure}

\begin{receta}[Diseño IMC]
\begin{enumerate}
  \item Obtenga un modelo nominal $\tilde G_p$ y factorícelo como en \eqref{eq:ps-factorizacion}.
  \item Forme $Q=(\tilde G_p^-)^{-1}F$ con un filtro que lo haga propio:
  \begin{itemize}
    \item cambios escalón de referencia: $F=1/(\lambda s+1)^n$;
    \item referencias tipo rampa: $F=(n\lambda s+1)/(\lambda s+1)^n$;
    \item rechazo de perturbaciones, procesos integradores o inestables: $F=(\gamma s+1)/(\lambda s+1)^n$, con $\gamma$ escogido para cancelar el polo más lento o inestable.
  \end{itemize}
  \item Ajuste $\lambda$ revisando la estabilidad con el proceso real o con los peores modelos posibles.
\end{enumerate}
\end{receta}

\begin{ejemplo}[label=ej:ps-bequette]{El modelo no coincide con el proceso}
El proceso real es $\Gp=1.5(1-s)/((s+1)(4s+1))$, pero el modelo es $\tilde G_p=2/(5s+1)$ (ejemplo de Bequette~\citep{Beq03}). ¿Para qué $\lambda$ es estable el lazo IMC?
\solucion Con $Q=(5s+1)/(2(\lambda s+1))$, la ecuación $1+Q(\Gp-\tilde G_p)=0$ da, multiplicando por el denominador común,
\[
  8\lambda s^3+(10\lambda-7.5)s^2+(2\lambda+6)s+1.5=0 .
\]
Routh exige $\lambda>0.75$ por el segundo coeficiente y, por la fila $s^1$, $(10\lambda-7.5)(2\lambda+6)>12\lambda$, es decir, $20\lambda^2+33\lambda-45>0$: \alerta{$\lambda>0.887$}. La \cref{fig:ps-robusto} lo confirma: con $\lambda=0.85$ el lazo es inestable; con $1.5$ responde rápido con algo de oscilación; con $3$, suave.
\end{ejemplo}

\begin{figure}[ht]
\centering
\includegraphics[width=0.85\linewidth]{cap08/f_imc_robusto}
\caption{Ejemplo~\ref{ej:ps-bequette}: el $\lambda$ del filtro decide si el error de modelado desestabiliza el lazo.}
\label{fig:ps-robusto}
\end{figure}

\codigorepo{cap08_pid_sintonizacion/C8_Ej5_IMC.m}{interna}{estabilidad interna con $Q_1$ y $Q_2$}
\codigorepo{cap08_pid_sintonizacion/C8_Ej5_IMC.m}{robusto}{ejemplo~\ref{ej:ps-bequette}}

\begin{modelo}[Modelo de Simulink: \texttt{crearModelo\_IMC}]
\texttt{codigo/cap08\_pid\_sintonizacion/crearModelo\_IMC.m} arma la \cref{fig:ps-imc} con el proceso y el modelo del ejemplo~\ref{ej:ps-bequette}, un escalón de referencia y una perturbación a la salida. Cambie \texttt{lambda} en el espacio de trabajo y vuelva a simular: por debajo de $0.887$ el lazo se inestabiliza.
\end{modelo}

\subsection{El PID que hay dentro del IMC}

Reordenando la \cref{fig:ps-imc}, el lazo IMC equivale a un lazo clásico con
\begin{equation}
  \Gc(s)=\frac{Q(s)}{1-\tilde G_p(s)Q(s)} .
  \label{eq:ps-imc-pid}
\end{equation}
Para muchos modelos sencillos, $\Gc$ resulta un PID (a veces con un filtro de primer orden): es la relación \clave{IMC-PID}~\citep{RMS86}.

\begin{receta}[IMC-PID]
\begin{enumerate}
  \item Forme $Q$ con un filtro que lo deje \emph{semipropio}: se admite que el numerador tenga un grado más, porque eso es justamente la acción derivativa. Para procesos integradores o inestables use $F=(\gamma s+1)/(\lambda s+1)^n$; si el modelo tiene un polo inestable en $s=p_u$, imponga $F(p_u)=1$.
  \item Calcule $\Gc$ con \eqref{eq:ps-imc-pid} y escríbalo como cociente de polinomios.
  \item Identifique $\Kc$, $\Ti$, $\Td$ y, si aparece, $\tau_f$ en $\Kc\Bigl(\dfrac{\Ti\Td s^2+\Ti s+1}{\Ti s}\Bigr)\dfrac{1}{\tau_f s+1}$.
  \item Ajuste $\lambda$ con simulaciones, con el modelo nominal y con modelos perturbados. Un punto de partida usual es entre un tercio y la mitad de la constante de tiempo dominante.
\end{enumerate}
\end{receta}

\paragraph{Primer orden.} Con $\tilde G_p=\Kproc/(\tau s+1)$ y $F=1/(\lambda s+1)$, $Q$ es \eqref{eq:ps-q-filtro} y
\[
  \Gc(s)=\frac{\tau s+1}{\Kproc\lambda s}:\qquad \Kc=\frac{\tau}{\Kproc\lambda},\quad\Ti=\tau ,
\]
el mismo PI de la síntesis directa: con modelo perfecto, el IMC es una síntesis directa con $T=F$. Queda un único parámetro; un $\lambda$ pequeño da un $\Kc$ grande, y al revés.

\paragraph{Segundo orden.} Con $\tilde G_p=\Kproc/((\tau_1s+1)(\tau_2s+1))$, $\Gc=(\tau_1s+1)(\tau_2s+1)/(\Kproc\lambda s)$: el PID ideal con $\Kc=(\tau_1+\tau_2)/(\Kproc\lambda)$, $\Ti=\tau_1+\tau_2$ y $\Td=\tau_1\tau_2/(\tau_1+\tau_2)$.

\paragraph{Primer orden con tiempo muerto.} El retardo no se invierte. Con la aproximación de Padé de primer orden de la \cref{sec:es-retardo}, $e^{-Ls}\approx(1-\frac{L}{2}s)/(1+\frac{L}{2}s)$, y la factorización pasa-todo,
\[
  \tilde G_p^+=\frac{1-\frac L2s}{1+\frac L2s},\qquad
  \tilde G_p^-=\frac{\Kproc}{(\tau s+1)(\frac L2s+1)},\qquad
  Q=\frac{(\tau s+1)(\frac L2s+1)}{\Kproc(\lambda s+1)} .
\]
Entonces $\Gc=\dfrac{(\tau s+1)(Ls+2)}{\Kproc(2\lambda+L)s}$, un PID ideal con
\begin{equation}
  \Kc=\frac{\tau+L/2}{\Kproc(\lambda+L/2)},\qquad \Ti=\tau+\frac L2,\qquad \Td=\frac{\tau L}{2\tau+L} .
  \label{eq:ps-imc-fopdt}
\end{equation}
Como Padé solo es buena a bajas frecuencias, $\lambda$ no puede ser arbitrariamente pequeño: Rivera, Morari y Skogestad recomiendan $\lambda>0.8L$~\citep{RMS86}.

\begin{ejemplo}[label=ej:ps-imc-fopdt]{IMC-PID con el retardo exacto}
Con $\Kproc=1.2$, $\tau=20$\,s y $L=4$\,s, evalúe $\lambda=3.2$, $6.7$ y $10$\,s simulando con el retardo exacto (estructura estándar, $N=10$).
\solucion $\Ti=22$\,s y $\Td=1.82$\,s en los tres casos; $\Kc=3.53$, $2.11$ y $1.53$. El sobrepico es $16$\,\%, $0$ y $0$, con $\ts=16.5$, $28.3$ y $41.1$\,s (\cref{fig:ps-imc-fopdt}). Bajando $\lambda$, el lazo con retardo exacto se inestabiliza hacia $\lambda\approx0.75$\,s: la regla $\lambda>0.8L=3.2$\,s deja un margen holgado.
\end{ejemplo}

\begin{figure}[ht]
\centering
\includegraphics[width=0.85\linewidth]{cap08/f_imc_fopdt}
\caption{Ejemplo~\ref{ej:ps-imc-fopdt}: IMC-PID sobre $1.2e^{-4s}/(20s+1)$ con el retardo exacto.}
\label{fig:ps-imc-fopdt}
\end{figure}

\codigorepo{cap08_pid_sintonizacion/C8_Ej5_IMC.m}{fopdt}{ejemplo~\ref{ej:ps-imc-fopdt}}

\paragraph{Proceso inestable.} Con $\tilde G_p=\Kproc/(1-\tau_u s)$, $\tau_u>0$, el polo está en $p_u=+1/\tau_u$. Un filtro $F=(\gamma s+1)/(\lambda s+1)^2$ con $F(1/\tau_u)=1$ exige
\[
  \frac{\gamma}{\tau_u}+1=\Bigl(\frac{\lambda}{\tau_u}+1\Bigr)^2
  \quad\Longrightarrow\quad
  \gamma=\lambda\Bigl(\frac{\lambda}{\tau_u}+2\Bigr),
\]
y con $Q=(1-\tau_us)(\gamma s+1)/(\Kproc(\lambda s+1)^2)$ resulta un PI:
\[
  \Kc=-\frac{\lambda+2\tau_u}{\Kproc\lambda},\qquad \Ti=\gamma .
\]
La ganancia es negativa porque, a alta frecuencia, $\Kproc/(1-\tau_us)\approx-\Kproc/(\tau_us)$ invierte el signo. Con $\Kproc=\tau_u=1$ los polos del lazo quedan dobles en $-1/\lambda$; para $\lambda=0.5$, $1$ y $2$\,s el sobrepico es $28$, $45$ y $79$\,\%: el cero del PI, en $-1/\gamma$, se acerca al origen a medida que $\lambda$ crece. Con un proceso inestable, el IMC se implementa como el lazo clásico \eqref{eq:ps-imc-pid}: la estructura de la \cref{fig:ps-imc} correría el modelo inestable en lazo abierto.

\begin{table}[ht]
\centering
\caption{Sintonías IMC-PID para el PID ideal (algunas con un filtro $1/(\tau_fs+1)$). Las filas con tiempo muerto usan aproximaciones: Taylor de primer orden, $e^{-Ls}\approx1-Ls$, para el PI, y Padé de primer orden para el PID; la fila con $\tau_3$ trunca además la serie de $1/(\tau_3s+1)$. Fuentes:~\citep{RMS86,Beq03}.}
\label{tab:ps-imc}
\renewcommand{\arraystretch}{1.7}
\begin{tabular}{@{}lccc@{}}
\toprule
Modelo $\tilde G_p$ & $\Kc\Kproc$ & $\Ti$ & $\Td$\\
\midrule
$\dfrac{\Kproc}{\tau s+1}$ & $\dfrac{\tau}{\lambda}$ & $\tau$ & ---\\
$\dfrac{\Kproc e^{-Ls}}{\tau s+1}$ (PI) & $\dfrac{\tau}{\lambda+L}$ & $\tau$ & ---\\
$\dfrac{\Kproc e^{-Ls}}{\tau s+1}$ (PID) & $\dfrac{\tau+L/2}{\lambda+L/2}$ & $\tau+\dfrac L2$ & $\dfrac{\tau L}{2\tau+L}$\\
$\dfrac{\Kproc}{(\tau_1s+1)(\tau_2s+1)}$ & $\dfrac{\tau_1+\tau_2}{\lambda}$ & $\tau_1+\tau_2$ & $\dfrac{\tau_1\tau_2}{\tau_1+\tau_2}$\\
$\dfrac{\Kproc(\tau_3s+1)e^{-Ls}}{(\tau_1s+1)(\tau_2s+1)}$ & $\dfrac{\tau_1+\tau_2-\tau_3}{\lambda+L}$ & $\tau_1+\tau_2-\tau_3$ & $\dfrac{\tau_1\tau_2}{\tau_1+\tau_2-\tau_3}-\tau_3$\\
$\dfrac{\Kproc}{\tau^2s^2+2\zeta\tau s+1}$ & $\dfrac{2\zeta\tau}{\lambda}$ & $2\zeta\tau$ & $\dfrac{\tau}{2\zeta}$\\
$\dfrac{\Kproc}{s}$ & $\dfrac{1}{\lambda}$ & --- & ---\\
$\dfrac{\Kproc e^{-Ls}}{s}$ & $\dfrac{2\lambda+L}{(\lambda+L)^2}$ & $2\lambda+L$ & ---\\
\bottomrule
\end{tabular}
\end{table}

%% ---------------------------------------------------------------------
\section{Desempeño contra robustez}
\label{sec:ps-robustez}
%% ---------------------------------------------------------------------

Una sintonía se calcula con un modelo, pero se aplica a un proceso que no es exactamente ese modelo: la ganancia cambia con el punto de operación, el tiempo muerto con el caudal, la constante de tiempo con el nivel. Una buena sintonía debe aguantar esas diferencias. Comparemos las reglas con el proceso del ejercicio de identificación, $1.2e^{-4s}/(20s+1)$, que tiene $L/\tau=0.2$.

\begin{ejemplo}[label=ej:ps-robustez]{Cuatro sintonías, nominal y perturbado}
Compare ZN, CC, ajustabilidad e IMC-PID ($\lambda=6.7$\,s) en estructura estándar con $N=10$: respuesta nominal, márgenes con el retardo exacto y estabilidad ante una malla de procesos con $\Kproc$ en $\pm30$\,\%, $\tau$ en $\pm30$\,\% y $L$ en $\pm50$\,\% (cinco valores de cada uno: 125 procesos).
\solucion Los márgenes de ganancia y de fase son los del capítulo~\ref{cap:rf} (definiciones~\ref{def:rf-gm} y \ref{def:rf-pm}); $\Delta L=\PM/\omega_{PM}$ es el retardo adicional que el lazo tolera antes de volverse inestable.
\begin{center}
\renewcommand{\arraystretch}{1.2}
\begin{tabular}{@{}lccccccc@{}}
\toprule
Sintonía & $\OS$ & $\ts$ (s) & IAE & GM & PM & $\Delta L$ (s) & estables\\
\midrule
Ziegler-Nichols & $87$\,\% & $42$ & $9.7$ & $1.33$ & $39^\circ$ & $2.2$ & $98/125$\\
Cohen-Coon & $95$\,\% & $48$ & $11.5$ & $1.31$ & $27^\circ$ & $1.3$ & $90/125$\\
Ajustabilidad & $0$\,\% & $26$ & $8.0$ & $2.91$ & $75^\circ$ & $10.1$ & $125/125$\\
IMC, $\lambda=6.7$ & $0$\,\% & $28$ & $8.7$ & $3.16$ & $76^\circ$ & $11.2$ & $125/125$\\
\bottomrule
\end{tabular}
\end{center}
Cohen-Coon, la más agresiva, tolera solo $1.3$\,s más de retardo (un $33$\,\% de $L$); en la esquina con $\Kproc$ alto, $\tau$ corto y $L$ largo, Ziegler-Nichols y Cohen-Coon se vuelven inestables. Las dos sintonías conservadoras toleran toda la malla. Para que Cohen-Coon la tolere hay que multiplicar su $\Kc$ por $0.48$; y entonces, curiosamente, \emph{mejora} también el desempeño nominal ($\OS=24$\,\%, IAE $9.4$), porque la sintonía original oscilaba demasiado.
\end{ejemplo}

\codigorepo{cap08_pid_sintonizacion/C8_Ej6_Robustez.m}{malla}{márgenes y malla de procesos con \cmd{allmargin}}

\begin{ideaclave}
\textbf{La agresividad se paga en robustez.}\enspace Las reglas que buscan la razón de un cuarto dejan márgenes de ganancia cercanos a $1.3$ y márgenes de fase de $25$ a $40^\circ$: muy poco para un proceso que cambia. Ante la perturbación, en cambio, son las mejores (IAE de $1.9$ contra $9.6$ y $10.4$ de las conservadoras). Escoger una sintonía es escoger un punto en ese compromiso, y el IMC lo hace explícito con un solo parámetro: $\lambda$.
\end{ideaclave}

%% ---------------------------------------------------------------------
\section{Para llevar}
\label{sec:ps-cierre}
%% ---------------------------------------------------------------------

\begin{parallevar}
\begin{enumerate}
  \item El controlador ve todo el lazo: válvula, proceso y transmisor. Se identifica moviendo la CO y midiendo la PV, con ganancias normalizadas a los rangos.
  \item Modelo FOPDT $\Kproc e^{-Ls}/(\tau s+1)$ con la tangente de máxima pendiente: $L$ en el corte con el valor inicial, $\tau$ hasta el $63$\,\% (o hasta el $100$\,\% en la variante de Ziegler y Nichols, que sobrestima $\tau$). En un primer orden la pendiente máxima está en $t=L$.
  \item Malla abierta: Ziegler-Nichols, Cohen-Coon (agresivas, buscan el decaimiento de un cuarto) y la regla empírica del coeficiente de ajustabilidad (conservadora). Todas dependen de $L/\tau$.
  \item Malla cerrada: $\Kcr$ y $\Pcr$ de la oscilación sostenida. Ziegler-Nichols (PI: $0.45\Kcr$, $\Pcr/1.2$) es oscilatoria; Tyreus-Luyben es más conservadora y \emph{no} es de un cuarto.
  \item Síntesis directa: $\Gc=\frac{1}{\Gp}\frac{T}{1-T}$. Con $T=1/(\lambda s+1)$, $\lambda$ es la constante de tiempo del lazo: menor $\lambda$, más rápido. Los ceros inestables y el retardo no se invierten: van en $T$.
  \item IMC: $Q=(\tilde G_p^-)^{-1}F$, estructura que retroalimenta solo el error del modelo, y $\Gc=Q/(1-\tilde G_pQ)$, que para modelos sencillos es un PID con un solo parámetro, $\lambda$. Revise siempre la estabilidad interna.
  \item Desempeño y robustez se oponen: las reglas agresivas dejan márgenes de ganancia cercanos a $1.3$ y se desestabilizan con cambios moderados del proceso.
\end{enumerate}
\end{parallevar}

\begin{profundiza}[Correcciones respecto a la versión anterior del capítulo]
(i) Se decía que en un primer orden la pendiente máxima coincide con el $63$\,\%; está al comienzo de la respuesta, en $t=L$. (ii) El PI de Ziegler-Nichols de malla cerrada es $0.45\Kcr$ y $\Pcr/1.2$, no $0.4\Kcr$ y $0.8\Pcr$. (iii) Tyreus-Luyben no es ``un método de caída de un cuarto'': es más conservador que Ziegler-Nichols a propósito. (iv) En síntesis directa se decía que con $\lambda$ pequeño ``el sistema no responde rápido''; es al revés: $\ts\approx4\lambda$. (v) En el ejemplo de Tyreus-Luyben el proceso aparecía como $\Gc$ y $\Pcr=6.2$ en lugar de $6.28$\,s; en el de fase no mínima, $\tau_2=2.5$ era $\Td=2.5$; en el IMC faltaba la $s$ en $U(s)=\frac{1}{\Kproc}(\tau s+1)R(s)$. (vi) Proceso inestable: la condición es $F(1/\tau_u)=1$, no $F(\tau_u)=1$. (vii) En la tabla IMC, la fila del primer orden con retardo mezclaba la ganancia del PID con el $\Ti$ del PI; se separó en una fila PI y otra PID, y se retiró la fila del subamortiguado con cero y retardo, cuyas fórmulas no se pudieron verificar. (viii) La ganancia de velocidad $K_v$ se define con la pendiente de la rampa. (ix) La tabla del coeficiente de ajustabilidad se conserva, marcada como regla empírica. (x) Se completaron las simulaciones pendientes y se calcularon todos los números. (xi) Notación: $\tau_i,\tau_d\to\Ti,\Td$; $\theta\to L$; $T$ del FOPDT $\to\tau$, para no confundirla con $T(s)$.
\end{profundiza}

%% ---------------------------------------------------------------------
\section{Problemas}
\label{sec:ps-problemas}
%% ---------------------------------------------------------------------

Los problemas marcados con \etiquetaM{} requieren \matlab/\simulink; los marcados con \etiquetaEstrella{} son de profundización.

\begin{table}[ht]
\centering
\caption{Respuesta experimental del proceso de temperatura del problema~\ref{prob:ps-ident}.}
\label{tab:ps-datos}
\small
\begin{tabular}{@{}l*{7}{c}@{}}
\toprule
$t$ (s) & 0 & 2 & 4 & 6 & 8 & 10 & 15\\
PV ($\grados$C) & 80.0 & 80.0 & 80.1 & 81.7 & 83.3 & 84.6 & 87.7\\
\midrule
$t$ (s) & 20 & 24 & 30 & 40 & 50 & 60 & 100\\
PV ($\grados$C) & 89.9 & 91.4 & 93.1 & 95.0 & 96.2 & 96.9 & 97.9\\
\bottomrule
\end{tabular}
\end{table}

\begin{problema}[identificación]\label{prob:ps-ident}
En un proceso de temperatura, la CO ($0$--$100$\,\%) pasa de $30$ a $45$\,\% y el transmisor está calibrado entre $50$ y $150\,\grados$C; los datos están en la \cref{tab:ps-datos}. (a) Halle $\Kproc$ normalizada con \eqref{eq:ps-ganancia} y en unidades físicas; ¿por qué coinciden aquí? (b) Con la tangente y el método del $63$\,\%, halle $L$ y $\tau$. (c) Repita con la variante de Ziegler y Nichols y compare $\tau$. (d) Calcule $r_c=L/\tau$ y ubique el proceso en la \cref{tab:ps-ajust}. (e) \etiquetaM{} Superponga el modelo $1.2e^{-4s}/(20s+1)$ sobre los datos y calcule el error cuadrático medio. (f) Derive $y(t)=\Kproc(1-e^{-t/\tau})$ y diga dónde está la pendiente máxima.
\end{problema}

\begin{problema}[reglas de malla abierta]\label{prob:ps-abierta}
Con $\Kproc=1.2$, $L=4$\,s y $\tau=20$\,s: (a) halle las sintonías P, PI y PID de Ziegler-Nichols; (b) las P, PI, PD y PID de Cohen-Coon; (c) la del coeficiente de ajustabilidad. (d) \etiquetaM{} Simule los tres PID con el retardo exacto (estructura estándar, $N=10$) y tabule $\OS$, $\ts$, IAE y $|u_c|_{\max}$ ante un escalón de referencia. (e) Ordénelos de más a menos agresivo. (f) \etiquetaM{} Repita (d) con un escalón de perturbación a la entrada del proceso. ¿Se mantiene el orden? ¿Por qué una regla puede ser mejor para seguir la referencia y peor para rechazar perturbaciones?
\end{problema}

\begin{problema}[malla cerrada]\label{prob:ps-cerrada}
Con $\Gp=1/((s+1)(2s+1)(4s+1))$ y control proporcional: (a) halle $\Kcr$, $\wcr$ y $\Pcr$ por sustitución $s=\jw$; (b) verifique con Routh y \etiquetaM{} simulando con $\Kc$ un poco por debajo, igual y un poco por encima de $\Kcr$; (c) sintonice un PID ideal con Ziegler-Nichols y con Tyreus-Luyben; (d) \etiquetaM{} simule y mida la relación de decaimiento; ¿cuál da el cuarto? (e) \etiquetaM{} multiplique la ganancia del proceso por $1.5$ y repita; ¿cuál tolera mejor el cambio?
\end{problema}

\begin{problema}[síntesis directa]\label{prob:ps-sintesis}
(a) Para $\Gp=2/(5s+1)$ y $T=1/(\lambda s+1)$, muestre que resulta un PI y halle $\Kc$ y $\Ti$. (b) Para $\Gp=3/((2s+1)(10s+1))$, muestre que resulta un PID ideal. (c) \etiquetaM{} Simule (b) con $\lambda=1$, $3$ y $5$\,s; reporte $\ts$ y $|u_c|_{\max}$ y enuncie la relación entre $\lambda$ y la velocidad. (d) Con $\Gp=(1-9s)/((15s+1)(3s+1))$, muestre que $T=1/(\lambda s+1)$ da un controlador inestable y que $T=(1-9s)/(\lambda s+1)^2$ no. (e) Lleve el último a la forma PID en serie con filtro y verifique para $\lambda=5$. (f) Explique por qué el cero inestable debe ir en $T$ y qué pasa con el grado relativo si no.
\end{problema}

\begin{problema}[diseño IMC]\label{prob:ps-imc}
(a) Para $\tilde G_p=2/(5s+1)$, forme $Q$ con un filtro de primer orden y verifique que con modelo perfecto $Y=FR$. (b) Halle $\Gc$ con \eqref{eq:ps-imc-pid} y compare con el problema~\ref{prob:ps-sintesis}(a). (c) Repita para $\Kproc/((\tau_1s+1)(\tau_2s+1))$ y verifique la fila de la \cref{tab:ps-imc}. (d) Para $\tilde G_p=1.5(1-s)/((s+1)(4s+1))$, factorice de forma simple y pasa-todo, forme los dos $Q$ y compare sus órdenes. (e) \etiquetaM{} Simule los dos diseños con $\lambda=2$ ante referencia y ante una perturbación $1/(4s+1)$ a la salida; ¿cuál prefiere? (f) ¿Qué característica de la respuesta queda fijada por el cero inestable, haga lo que haga el controlador?
\end{problema}

\begin{problema}[robustez]\label{prob:ps-robustez}
Con las tres sintonías PID del problema~\ref{prob:ps-abierta}: (a) \etiquetaM{} calcule los márgenes de ganancia y de fase con el retardo exacto; (b) halle el retardo adicional $\Delta L$ que tolera cada una, en segundos y en porcentaje de $L$; (c) \etiquetaM{} varíe $\Kproc$ en $\pm30$\,\%, $\tau$ en $\pm30$\,\% y $L$ en $\pm50$\,\% (cinco valores de cada uno) y cuente los lazos estables; (d) identifique el peor caso y confírmelo por simulación; (e) ordene las sintonías por robustez y compare con el orden por agresividad; (f) \etiquetaEstrella{} halle el factor por el que hay que multiplicar el $\Kc$ de la menos robusta para que tolere toda la malla, y diga cuánto desempeño nominal cuesta.
\end{problema}

\begin{problema}[IMC-PID con tiempo muerto]\label{prob:ps-imc-fopdt}
Con $\Kproc=1.2$, $\tau=20$\,s y $L=4$\,s: (a) aproxime el retardo por Padé de primer orden, factorice y forme $Q$; (b) obtenga $\Gc$ y deduzca \eqref{eq:ps-imc-fopdt}; (c) compare con la fila FOPDT-PI de la \cref{tab:ps-imc}: ¿de qué aproximación sale cada una? (d) \etiquetaM{} Evalúe $\lambda=3.2$, $6.7$ y $10$\,s con el retardo exacto y tabule $\OS$, $\ts$, IAE y $|u_c|_{\max}$. (e) \etiquetaM{} Halle el menor $\lambda$ estable y compárelo con $0.8L$. (f) Compare la mejor sintonía IMC con las del problema~\ref{prob:ps-abierta} en IAE y margen de fase.
\end{problema}

\begin{problema}[estabilidad interna y errores de modelado]\label{prob:ps-interna}
(a) Para el ejemplo~\ref{ej:ps-bequette}, deduzca la ecuación característica y la cota de $\lambda$ por Routh. (b) \etiquetaM{} Simule por debajo, en y por encima de la cota. (c) \etiquetaM{} Con $\Gp=(1-9s)/((15s+1)(3s+1))$ y $\lambda=5$, simule $Q_1$ y $Q_2$ y grafique $y$ y $u$ durante $200$\,s. (d) Halle el polo responsable de la divergencia, su constante de tiempo y el instante en que $|u|$ supera $50$; explique por qué $y$ no lo revela. (e) Defina estabilidad interna. (f) Si el proceso es $K/(5s+1)$ y el modelo $2/(5s+1)$, halle la ecuación característica del lazo IMC y el rango de $K$ estable. ¿Cuánto error de ganancia tolera?
\end{problema}

%% ---------------------------------------------------------------------
\section{Soluciones seleccionadas}
%% ---------------------------------------------------------------------

\paragraph{Problema~\ref{prob:ps-ident}}
(a) $\Delta y\approx17.9\,\grados$C: $\Kproc=(17.9/100)/(15/100)=1.19\approx1.2$, y en unidades físicas $1.19\,\grados$C/\%; coinciden porque el rango de la PV es de $100\,\grados$C. (b) La pendiente máxima de los datos, $0.8\,\grados$C/s entre $4$ y $6$\,s, da $L=3.9$\,s; el $63.2$\,\% ($91.3\,\grados$C) se alcanza en $t\approx23.8$\,s, así que $\tau_{63}\approx19.9$\,s. (c) $\tau_{ZN}=17.9/0.8=22.4$\,s, mayor. (d) $r_c\approx0.2$: justo en el límite entre la fila PI y la PID; se usa la PID. (e) ECM $=0.0019\,(\grados\text{C})^2$. (f) La pendiente $\frac{\Kproc}{\tau}e^{-t/\tau}$ es máxima en $t=0$; en $t=\tau$ vale el $37$\,\% de ese máximo.

\paragraph{Problema~\ref{prob:ps-abierta}}
(a) P: $\Kc=4.17$; PI: $3.75$, $13.3$\,s; PID: $5$, $8$\,s, $2$\,s. (b) P: $4.44$; PI: $3.82$, $9.42$\,s; PD: $5.35$, $\Td=0.99$\,s; PID: $5.76$, $9.10$\,s, $1.40$\,s. (c) $2.29$, $22$\,s, $1.82$\,s. (d) ZN: $\OS=87$\,\%, $\ts=42$\,s, IAE $9.7$, $|u_c|_{\max}=54$; CC: $95$\,\%, $48$\,s, $11.5$, $61$; ajustabilidad: $0$\,\%, $26$\,s, $8.0$, $25$. (e) CC, ZN, ajustabilidad. (f) Picos $0.229$, $0.231$ y $0.288$; IAE $1.90$, $1.89$ y $9.6$: el orden se invierte.

\paragraph{Problema~\ref{prob:ps-cerrada}}
(a) $8s^3+14s^2+7s+1+\Kc=0$: $\wcr=\sqrt{7/8}=0.935$\,rad/s, $\Kcr=11.25$, $\Pcr=6.72$\,s. (c) ZN: $6.75$, $3.36$\,s, $0.84$\,s; TL: $5.11$, $14.8$\,s, $1.07$\,s. (d) ZN: $\OS=48$\,\%, decaimiento $0.16$; TL: $\OS=11$\,\%, sin segundo pico. (e) Con la ganancia $\times1.5$: $56$\,\% y $23$\,\%.

\paragraph{Problema~\ref{prob:ps-sintesis}}
(a) $\Kc=5/(2\lambda)$, $\Ti=5$\,s. (b) $\Kc=4/\lambda$, $\Ti=12$\,s, $\Td=5/3$\,s. (c) $\ts=3.9$, $12.1$ y $19.8$\,s; $|u_c|_{\max}=43.8$, $14.6$ y $8.8$: menor $\lambda$, más rápido. (d) Polo en $+1/9$. (e) $\Kc=18/(2\lambda+9)$, $\Ti=18$\,s, $\Td=2.5$\,s, $\tau_f=\lambda^2/(2\lambda+9)$; con $\lambda=5$, $0.947$ y $1.32$\,s.

\paragraph{Problema~\ref{prob:ps-imc}}
(d) Simple: $\tilde G_p^+=1-s$, $Q=(s+1)(4s+1)/(1.5(\lambda s+1)^2)$, de orden $2$. Pasa-todo: $\tilde G_p^+=(1-s)/(1+s)$, $\tilde G_p^-=1.5/(4s+1)$, $Q=(4s+1)/(1.5(\lambda s+1))$, de orden $1$. (e) Con $\lambda=2$: $\ts=12.5$ y $10.0$\,s, IAE $5$ y $4$ ($2\lambda+1$ y $\lambda+2$); ante la perturbación, picos $0.52$ y $0.47$. Se prefiere la pasa-todo. (f) La respuesta inversa.

\paragraph{Problema~\ref{prob:ps-robustez}}
(a) GM $1.33$, $1.31$, $2.91$; PM $39^\circ$, $27^\circ$, $75^\circ$. (b) $\Delta L=2.2$\,s ($54$\,\%), $1.3$\,s ($33$\,\%) y $10.1$\,s ($252$\,\%). (c) $98$, $90$ y $125$ de $125$. (d) $1.3\Kproc$, $0.7\tau$, $1.5L$. (f) CC con $\Kc\times0.48=2.77$: $\OS=24$\,\%, $\ts=31$\,s, IAE $9.4$, mejor que la original.

\paragraph{Problema~\ref{prob:ps-imc-fopdt}}
(c) La fila PI sale de $e^{-Ls}\approx1-Ls$; la PID, de Padé. La versión anterior de la tabla combinaba la ganancia de la segunda con el $\Ti$ de la primera. (d) $\Kc=3.53$, $2.11$, $1.53$; $\OS=16$\,\%, $0$, $0$; $\ts=16.5$, $28.3$, $41.1$\,s; IAE $5.8$, $8.7$, $12.0$; $|u_c|_{\max}=38$, $23$, $16$. (e) $\lambda_{\min}\approx0.75$\,s frente a $3.2$\,s. (f) IMC con $\lambda=6.7$: IAE $8.7$ y PM $76^\circ$, frente a $9.7$ y $39^\circ$ de Ziegler-Nichols.

\paragraph{Problema~\ref{prob:ps-interna}}
(a) $\lambda>0.887$. (d) Polo en $+1/9$, constante de tiempo $9$\,s; $|u|>50$ en $t=27.95$\,s. El cero del proceso cancela el polo en la función $Y/R$. (f) $1+Q(\Gp-\tilde G_p)=0$ se reduce a $2\lambda s+K=0$: estable para todo $K>0$. Con un error puro de ganancia, el IMC con modelo de primer orden tolera cualquier ganancia positiva.

\respuestaspausas

